% META: {"id": "TSPE-CONTIN-EX-035", "chapitre": "TSPE-CONTINUITE", "type_objet": "exercice", "capacites": ["TSPE-CONTIN-C2"], "capacites_codes": ["C2"], "methodes": ["M2"], "parcours": 1, "competences": ["chercher", "calculer", "raisonner"], "duree_min": 12, "mode_creation": "ex_nihilo", "sources_inspiration": [], "fichier_tex": "chapitres/TSPE-CONTINUITE/exercices/TSPE-CONTIN-EX-035.tex", "corrige_tex": "chapitres/TSPE-CONTINUITE/corriges/TSPE-CONTIN-CO-035.tex", "status": "generated"}
% BEGIN-VERIFY
% from sympy import symbols, diff, Rational, sqrt, simplify, solveset, S as SS, Eq, nsolve
% x = symbols('x')
% def iterer(f, u0, n):
%     u = u0
%     out = [u]
%     for _ in range(n):
%         u = f.subs(x, u)
%         out.append(u)
%     return out
% f = (x + 3)/2
% assert solveset(Eq(f, x), x, SS.Reals) == {3}
% t = [float(v) for v in iterer(f, 1, 4)]
% assert t == [1.0, 2.0, 2.5, 2.75, 2.875]
% END-VERIFY
\begin{exercice}{TSPE-CONTIN-EX-035}{1}{12}
On considere la suite $(v_n)$ definie par $v_0 = 1$ et $v_{n+1} = \dfrac{v_n + 3}{2}$.
\begin{enumerate}
  \item Calculer $v_1$, $v_2$, $v_3$ et $v_4$.
  \item Montrer que l'intervalle $[1\,;3]$ est stable par $f : x \mapsto \dfrac{x+3}{2}$, puis que $1 \leq v_n \leq 3$ pour tout $n$.
  \item Montrer que $(v_n)$ est croissante.
  \item En deduire que $(v_n)$ converge et determiner sa limite.
  \item On pose $w_n = 3 - v_n$. Montrer que $(w_n)$ est geometrique, en preciser la raison, puis retrouver la limite de $(v_n)$ par un calcul direct.
\end{enumerate}
\end{exercice}
